\scanpage{59}
Ce qui donne des conditions suffisantes pour la validité de c). Cependant, elles ne sont pas nécessaires [\unsure{$=$|$\bar m$} en exigeant la validité de c pour $i<n$, et tous les tordus de $F$], comme on voit en faisant $n=1$ : alors il suffit qu'on ait $\operatorname{prof} F_x\geq2$ pour $x$ fermé dans $U$, pour pouvoir conclure que
\[
H^0(U,F)\mathrel{\unsure{-|\to}}H^0(\widehat X,\widehat F)
\]
est bijectif.

\textcircled{\scriptsize 2\textsuperscript{o}} La suite exacte nous apprend que \unsure{l'ar.|l'hom.|l'homom.} $H^i(U,F)\to H^{i+1}_Z(X,F)$ est un isom mod vectoriels finis. Donc
\[
\deleted{H^i(U,F}
\]
\margin{\textcircled{$A'_i$}}
\[
\left\{\begin{aligned}
&H^i(U,F)\text{ de dim finie sur }k\\
&\Updownarrow\\
&H^{i+1}_z(F_z)\text{. de dim finie [i.e. de \unsure{long.|lg.} finie] pour tt }z\in Z\\
&\Updownarrow\\
&z\deleted{\unsure{\notin|\subset}}\in Z\ \Longrightarrow\ z\text{ pt isolé de }\operatorname{supp} E^{r-i-1}\cup\{z\}\\
&\Updownarrow\\
&\alpha_{i+1}\text{ injectif pour }m\text{ grand}\\
&\Updownarrow\\
&H^i(X,F(-m))\mathrel{\unsure{-|\to}}H^i(\deleted{X}U,F(-m))\ \text{surjectif pour }m\text{ grand}
\end{aligned}\right.
\]
% archive sheet number: 58
